\documentclass[11pt]{amsart}
\newtheorem{lemma}{Lemma}
\begin{document}

\title{Answer to Question \#9}

\author{From work by D.~Miao, G.~Lerman, J.~Kileel}

\maketitle

Yes, such algebraic relations do exist.  
Assemble the various tensors $\{Q^{(\alpha \beta \gamma \delta)} : \alpha, \beta, \gamma, \delta \in [n] \}$ into one tensor $\mathbf{Q} \in \mathbb{R}^{3n \times 3n \times 3n \times 3n}$, thought of as an $n \times n \times n \times n$ block tensor where the $(\alpha, \beta, \gamma, \delta)$-block is $Q^{(\alpha \beta \gamma \delta)} \in \mathbb{R}^{3 \times 3 \times 3 \times 3}$. 
Let $\mathbf{F}$ be the polynomial map sending $\{Q^{(\alpha \beta \gamma \delta)} : \alpha, \beta, \gamma, \delta \in [n] \}$ to the $5 \times 5$ minors of the four $3n \times 27n^3$ matrix flattenings of $\mathbf{Q}$.
We will prove that $\mathbf{F}$ satisfies the desired properties.

A key point is to discover the following algebraic identity.  


\begin{lemma}
\label{lem:identity}
Consider $\mathbf{Q} \in \mathbb{R}^{3n \times 3n \times 3n \times 3n}$ as above.  
It admits a Tucker tensor decomposition
\begin{equation}\label{eq:tucker}
\mathbf{Q} = \mathcal{C} \times_1 \mathbf{A} \times_2 \mathbf{A} \times_3 \mathbf{A} \times_4 \mathbf{A},
\end{equation}
for $\mathcal{C} \in \mathbb{R}^{4 \times 4 \times 4 \times 4}$ and $\mathbf{A} \in \mathbb{R}^{3n \times 4}$.  
Explicitly, we can take 
\begin{equation*}
\mathcal{C}_{abcd} = 
\begin{cases}
\operatorname{sgn}(abcd) & \text{ if } a,b,c,d \in [4] \text{ are distinct} \\
0 & \text{ otherwise},
\end{cases}
\end{equation*}
where sgn is parity of a permutation, and $\mathbf{A}$ to be the vertical concatenation $[A^{(1)}; \dots ; A^{(n)}]$.
\end{lemma}
\begin{proof}
Let $[n] \times [3]$ stand for the indices of $\mathbf{Q}$ in each mode and for the row indices of $\mathbf{A}$.
By definition of Tucker product, for all $(\alpha, i), (\beta, j), (\gamma, k), (\delta, \ell) \in [n] \times [3]$ we have
\begin{align*}
  &\left(\mathcal{C} \times_1 \mathbf{A} \times_2 \mathbf{A} \times_3 \mathbf{A} \times_4 \mathbf{A}\right)_{(\alpha, i), (\beta, j), (\gamma, k), (\delta, \ell)}  \\[0.25em] 
  &= \sum_{a,b,c,d \in [4]} \mathcal{C}_{abcd} \mathbf{A}_{(\alpha, i), a} \mathbf{A}_{(\beta, j), b} \mathbf{A}_{(\gamma, k), c} \mathbf{A}_{(\delta, \ell), d} \\[0.25em]
  &= \sum_{a,b,c,d \in [4] \text{ distinct}} \operatorname{sgn}(abcd) A^{(\alpha)}_{ia} A^{(\beta)}_{jb} A^{(\alpha)}_{kc} A^{(\alpha)}_{\ell d} \\[0.25em]
  &= \det [A^{(\alpha)}(i, :); A^{(\beta)}(j, :); A^{(\gamma)}(k, :); A^{(\delta)}(\ell, :)] \\[0.25em]
  &= Q^{(\alpha \beta \gamma \delta)}_{i j k \ell} = \mathbf{Q}_{(\alpha, i), (\beta, j), (\gamma, k), (\delta, \ell)}. \qedhere
\end{align*} 
\end{proof}


The lemma explains why $\mathbf{F}$ captures algebraic relations between the tensors $\{Q^{(\alpha \beta \gamma \delta)} : \alpha, \beta, \gamma, \delta \in [n] \}$.  
Indeed, the block tensor $\mathbf{Q}$ has multilinear rank bounded by $(4,4,4,4)$ due to the Tucker decomposition in \eqref{eq:tucker}.
Therefore, all $5 \times 5$ minors in $\mathbf{F}$ vanish.


Below we break up the proof of the third property into two directions.  
The other properties are clear.   
Throughout the proof, for $\lambda \in \mathbb{R}^{n \times n \times n \times n}$ we let $\lambda \odot_{b} \mathbf{Q} \in \mathbb{R}^{3n \times 3n \times 3n \times 3n}$ denote blockwise scalar multiplication, i.e., 
the $(\alpha, \beta, \gamma, \delta)$-block of $\lambda \odot_{b} \mathbf{Q}$ is $\lambda_{\alpha \beta \gamma \delta} Q^{(\alpha \beta \gamma \delta)} \in \mathbb{R}^{3 \times 3 \times 3 \times 3}$.
Roughly speaking, we need to show that a blockwise scaling of $\mathbf{Q}$ preserves multilinear rank if and only if the scaling is a rank-1 tensor off the diagonal.

\medskip

\subsection*{``If" Direction}
This follows easily from Lemma~\ref{lem:identity}.
Assume $\lambda \in \mathbb{R}^{n \times n \times n \times n}$ agrees off-diagonal with $u \otimes v \otimes w \otimes x$ for $u,v,w,x \in (\mathbb{R}^{*})^n$ and is $0$ on the diagonal.
Then
\begin{equation*}
\lambda \odot_b \mathbf{Q} = (u \otimes v \otimes w \otimes x) \odot_b \mathbf{Q},
\end{equation*}
because the diagonal blocks of $\mathbf{Q}$ vanish.  That is,
$
Q^{(\alpha \alpha \alpha \alpha)} = 0
$
since each entry of $Q^{(\alpha \alpha \alpha \alpha)}$ is the determinant of a matrix with a repeated row.
Note that blockwise scalar product with a rank-1 tensor with nonzero entries is equivalent to Tucker product with invertible matrices:
\begin{equation*}
(u \otimes v \otimes w \otimes w) \odot_b \mathbf{Q} = \mathbf{Q} \times_1 D_u \times_2 D_v \times_3 D_w \times_4 D_x.
\end{equation*}
Here $D_u \in \mathbb{R}^{3n \times 3n}$ is the diagonal matrix triplicating the entries of $u$ and likewise for $D_v, D_w, D_x$.
Thus $\lambda \odot_b \mathbf{Q}$ and $\mathbf{Q}$ have the same multilinear rank, and from the lemma $\mathbf{F}(\lambda_{\alpha \beta \gamma \delta} Q^{(\alpha \beta \gamma \delta)} : \alpha, \beta, \gamma, \delta \in [n]) = 0$.

\medskip

\subsection*{``Only If" Direction}

The converse takes more work. 
Let $\lambda \in \mathbb{R}^{n \times n \times n \times n}$ have nonzero entries precisely off the diagonal and assume $\mathbf{F}(\lambda_{\alpha \beta \gamma \delta} Q^{(\alpha \beta \gamma \delta)} : \alpha, \beta, \gamma, \delta \in [n]) = 0$.
We further assume $\lambda_{\alpha 1 1 1 } = \lambda_{1 \beta 1 1} = \lambda_{1 1 \gamma 1} = \lambda_{1 1 1 \delta}=1$ for all $\alpha, \beta, \gamma, \delta \in \{2, \ldots, n\}$.
We reduce to this case by replacing $\lambda$ by its entrywise product with $\bar{u} \otimes \bar{v} \otimes \bar{w} \otimes \bar{x}$, where 
\begin{equation*}
\bar{u}_{\alpha} = \begin{cases} 
1 & \text{for } \alpha = 1\\
\lambda_{\alpha 1 1 1}^{-1} & \text{for } \alpha \in \{2, \dots, n\}, 
\end{cases}
\end{equation*}
and $\bar{v}, \bar{w}, \bar{x}$ are defined similarly using the second, third and fourth modes respectively. 
The replacement preserves the multilinear rank of $\lambda \odot_b \mathbf{Q}$ and whether or not $\lambda$ agrees off-diagonal with a rank-$1$ tensor.  Hence it is without loss of generality.


Through some explicit calculations, we will prove there exists $c \in \mathbb{R}^*$ such that
\begin{itemize}
    \item $\lambda_{\alpha \beta \gamma \delta} = c$ if exactly two of $\alpha, \beta, \gamma, \delta$ equal $1$
    \item $\lambda_{\alpha \beta \gamma \delta} = c^2$ if exactly one of $\alpha, \beta, \gamma, \delta$ equals $1$
    \item $\lambda_{\alpha \beta \gamma \delta} = c^3$ if none of $\alpha, \beta, \gamma, \delta$ equal $1$ and $\alpha, \beta, \gamma, \delta$ are not identical.
\end{itemize}
This will establish the ``only if" direction, as setting $u = v= w= (1, c, \ldots, c)$ and $x= (\tfrac{1}{c}, 1, \ldots, 1)$ gives $\lambda_{\alpha \beta \gamma \delta} = u_{\alpha} v_{\beta} w_{\gamma} x_{\delta}$ whenever $\alpha, \beta, \gamma, \delta$ are not identical.
Our proof strategy is to examine appropriate coordinates 
of $\mathbf{F}(\lambda_{\alpha \beta \gamma \delta} Q^{(\alpha \beta \gamma \delta)} : \alpha, \beta, \gamma, \delta \in [n]) = 0$ in order to constrain $\lambda$.
Equivalently, we will consider the vanishing of the determinants of certain well-chosen $5 \times 5$ submatrices of the flattenings of $\lambda \odot_b \mathbf{Q}$.  
Write $\mathbf{Q}_{(1)}$ and $(\lambda \odot_b \mathbf{Q})_{(1)}$ for mode-1 flattenings in $\mathbb{R}^{3n \times 27n^3}$.  
Rows correspond to the first tensor mode and are indexed by $(\alpha, i) \in [n] \times [3]$, while columns correspond to the other  modes and are indexed by $((\beta, j), (\gamma, k), (\delta, \ell)) \in ([n] \times [3])^{3}$.

\medskip

\underline{{Step 1}}:
The first submatrix of $(\lambda \odot_b \mathbf{Q})_{(1)}$ we consider has column indices
$((\alpha,1),(1,3),(1,2))$, 
$((1,2),(\beta,2),(1,1))$,
$((1,2),(\beta,3),(1,1))$,
$((1,3),(\beta,3),$ $(1,2))$,
$((1,1),(\beta,1),(1,3))$
and row indices 
$(1,1)$,
$(1,2)$,
$(1,3)$,
$(\alpha,1)$,
$(\alpha,2)$, 
where $\alpha, \beta \in \{2, \ldots, n\}$.
Explicitly, the submatrix \nolinebreak is  
\begin{equation*}
\begin{bmatrix}
{Q}^{(1\alpha11)}_{1132} & {Q}^{(11\beta1)}_{1221} & {Q}^{(11\beta1)}_{1231} & {Q}^{(11\beta1)}_{1332} & {Q}^{(11\beta1)}_{1113} \\[5pt] 
{Q}^{(1\alpha11)}_{2132} & {Q}^{(11\beta1)}_{2221} & {Q}^{(11\beta1)}_{2231} & {Q}^{(11\beta1)}_{2332} & {Q}^{(11\beta1)}_{2113} \\[5pt]
{Q}^{(1\alpha11)}_{3132} & {Q}^{(11\beta1)}_{3221} & {Q}^{(11\beta1)}_{3231} & {Q}^{(11\beta1)}_{3332} & {Q}^{(11\beta1)}_{3113} \\[4pt]
\lambda_{\alpha\alpha11}{Q}^{(\alpha\alpha11)}_{1132} & \lambda_{\alpha1\beta1}{Q}^{(\alpha1\beta1)}_{1221} & \lambda_{\alpha1\beta1}{Q}^{(\alpha1\beta1)}_{1231} & \lambda_{\alpha1\beta1}{Q}^{(\alpha1\beta1)}_{1332} & \lambda_{\alpha1\beta1}{Q}^{(\alpha1\beta1)}_{1113} \\[4pt]
\lambda_{\alpha\alpha11}{Q}^{(\alpha\alpha11)}_{2132} & \lambda_{\alpha1\beta1}{Q}^{(\alpha1\beta1)}_{2221} & \lambda_{\alpha1\beta1}{Q}^{(\alpha1\beta1)}_{2231} & \lambda_{\alpha1\beta1}{Q}^{(\alpha1\beta1)}_{2332} & \lambda_{\alpha1\beta1}{Q}^{(\alpha1\beta1)}_{2113}
\end{bmatrix},
\end{equation*}
which we abbreviate as 
\begin{equation} \label{eq:my-det1}
\begin{bmatrix}
\ast & \ast & \ast & \ast & \ast \\
\ast & \ast & \ast & \ast & \ast \\
\ast & \ast & \ast & \ast & \ast \\
\lambda_{\alpha\alpha11}\ast & \lambda_{\alpha1\beta1}\ast & \lambda_{\alpha1\beta1}\ast & \lambda_{\alpha1\beta1}\ast & \lambda_{\alpha1\beta1}\ast \\
\lambda_{\alpha\alpha11}\ast & \lambda_{\alpha1\beta1}\ast & \lambda_{\alpha1\beta1}\ast & \lambda_{\alpha1\beta1}\ast & \lambda_{\alpha1\beta1}\ast \\
\end{bmatrix},
\end{equation}
with asterisk denoting the corresponding entry in $\mathbf{Q}_{(1)}$. 
We view the determinant of \eqref{eq:my-det1} as a polynomial with respect to $\lambda$.
It has degree $\leq 2$ in the variables $\lambda_{\alpha \alpha 11}, \lambda_{\alpha 1 \beta 1}$. 
Observe that if $\lambda_{\alpha 1 \beta 1} = 0$, the bottom two rows of the matrix are linearly independent.  
Also if $\lambda_{\alpha 1 \beta 1} - \lambda_{\alpha \alpha 11} = 0$, then 
\eqref{eq:my-det1} equals a $5 \times 5$ submatrix of $\mathbf{Q}_{(1)}$ with rows operations performed; therefore \eqref{eq:my-det1} is rank-deficient.
It follows that the determinant of \eqref{eq:my-det1} takes the form
\begin{equation*}
s \lambda_{\alpha 1 \beta 1}(\lambda_{\alpha 1 \beta 1} - \lambda_{\alpha \alpha 1 1}).
\end{equation*}
Here the scale $s = s(A^{(1)}, A^{(\alpha)}, A^{(\beta)})$ is a polynomial in the $A$-matrices. 
Due to polynomiality, $s$ is nonzero Zariski-generically if we can exhibit a \textit{single} instance of matrices $A^{(1)}, A^{(\alpha)}, A^{(\beta)}$ where the determinant of \eqref{eq:my-det1} does not vanish identically for all $\lambda_{\alpha 1 \beta 1}, \lambda_{\alpha \alpha 1 1}$.
Furthermore, we just need an instance with $\alpha = \beta$, as this corresponds to a specialization of the case $\alpha \neq \beta$.  
Computational verification with a random numerical instance of $A^{(1)}, A^{(\alpha)}$ proves the non-vanishing (see attached code).  
Recalling the standing assumptions, we deduce 
$\lambda_{\alpha 1 \beta 1} = \lambda_{\alpha \alpha 1 1}$.


We apply the same argument to modewise permutations of $\lambda \odot_b \mathbf{Q}$ and $\mathbf{Q}$, and obtain
\begin{equation*}\label{eq:nice-1}
\lambda_{\pi(\alpha 1 \beta 1)} = \lambda_{\pi(\alpha \alpha 1 1)} \quad \text{for all } \alpha, \beta \in \{2, \dots, n\} \text{ and permutations } \pi.
\end{equation*}
The argument goes through as $\pi \cdot \mathbf{Q}$ and $\pi \cdot (\lambda \odot_b \mathbf{Q})$ have multilinear ranks bounded by $(4,4,4,4)$ and $\pi \cdot \mathbf{Q} = \operatorname{sgn}(\pi) \mathbf{Q}$.  
So \eqref{eq:my-det1} looks the same but with indices permuted and possibly a sign flip.

 
We now see that $\lambda$-entries with  two $1$-indices agree.
Indeed, taking $\alpha = \beta$ above gives $\lambda_{\pi_1(\alpha 1 \alpha 1)} = \lambda_{\pi_2(\alpha \alpha 1 1)}$ for all $\pi_1$ and $\pi_2$ that fix $(\alpha \alpha 1 1)$ and $(\alpha 1 \alpha 1)$ respectively.  
So, $\lambda_{\alpha \alpha 1 1} = \lambda_{\pi(\alpha \alpha 1 1)}$ for all $\pi$.  
Taking $\alpha \neq \beta$ gives $\lambda_{\alpha \alpha 1 1} = \lambda_{\pi(\alpha 1 \beta 1)} = \lambda_{\beta \beta 1 1}$ for all $\pi$.  
Together, there exists $c \in \mathbb{R}^*$ such that 
$c = \lambda_{\pi(\alpha \beta 1 1)}$ for all $\alpha, \beta \in \{2, \dots, n\}$ and permutations \nolinebreak $\pi$.

\medskip 

\underline{Step 2}:  
Next we consider the submatrix of $(\lambda \odot_b \mathbf{Q})_{(1)}$ with column indices  
$((\beta,1),(\gamma,3),(1,2))$, 
$((1,2),(\beta,2),(1,1))$,
$((1,2),(\beta,3),(1,1))$,
$((1,3),(\beta,3),$ $(1,2))$,
$((1,1),(\beta,1),(1,3))$
and row indices 
$(1,1)$,
$(1,2)$,
$(1,3)$,
$(\alpha,1)$,
$(\alpha,2)$, 
where $\alpha, \beta, \gamma \in \{2, \ldots, n\}$.
It looks like 
\begin{equation} \label{eq:my-det2}
\begin{bmatrix}
c\ast & \ast & \ast & \ast & \ast \\
c\ast & \ast & \ast & \ast & \ast \\
c\ast & \ast & \ast & \ast & \ast \\
\lambda_{\alpha \beta \gamma 1}\ast & c\ast & c\ast & c\ast & c\ast \\
\lambda_{\alpha \beta \gamma 1}\ast & c\ast & c\ast & c\ast & c\ast \\
\end{bmatrix},
\end{equation}
where asterisks denote corresponding entries in $\mathbf{Q}_{(1)}$.
As a polynomial in $c$ and $\lambda_{\alpha \beta \gamma 1}$, the determinant of \eqref{eq:my-det2} is a scalar multiple of $c(c^2 - \lambda_{\alpha \beta \gamma 1})$.
This is because the polynomial has degree $\leq 3$, if $c=0$ then the bottom two rows of \eqref{eq:my-det2} are linearly dependent, and if $c^2 = \lambda_{\alpha \beta \gamma 1}$ then \eqref{eq:my-det2} is a $5 \times 5$ submatrix of $\mathbf{Q}_{(1)}$ with row and column operations performed.  
The scale is a polynomial in $A^{(1)}, A^{(\alpha)}, A^{(\beta)}, A^{(\gamma)}$.
It is Zariski-generically nonzero if we exhibit one instance of $A$-matrices such that the determinant of \eqref{eq:my-det1} does not vanish for all $c, \lambda_{\alpha \beta \gamma 1}$.
Further, it suffices to find an instance where $\alpha = \beta = \gamma$, as all other cases specialize to this.  Computational verification with a random numerical instance of $A^{(1)}, A^{(\alpha)}$ proves the non-vanishing.  
It follows that $c^2 = \lambda_{\alpha \beta \gamma 1}$. 
Appealing to symmetry like before, 
$
c^2 = \lambda_{\pi(\alpha \beta \gamma 1)}
$
for all $\alpha, \beta, \gamma \in \{2, \dots, n\}$ and permutations $\pi$.  
Summarizing, all $\lambda$-entries with a single $1$-index equal $c^2$.

\medskip

\underline{Step 3}:
Consider  the submatrix of $(\lambda \odot \mathbf{Q})_{(1)}$ with columns 
$((\beta,1),(\gamma,3),$ $(\delta, 2))$, 
$((1,2),(\alpha,2),$ $(1,1))$,
$((1,2),(\alpha,3),(1,1))$,
$((1,3),(\alpha,3),(1,2))$,
$((1,1),$ $(\alpha,1),(1,3))$
and rows 
$(1,1)$,
$(1,2)$,
$(1,3)$,
$(\alpha,1)$,
$(\alpha,2)$, 
where $\alpha, \beta, \gamma, \delta \in \{2, \ldots, n\}$ and $\alpha, \delta$ are distinct.
The submatrix looks like 
\begin{equation} \label{eq:my-det3}
\begin{bmatrix}
c^2\ast & \ast & \ast & \ast & \ast \\
c^2\ast & \ast & \ast & \ast & \ast \\
c^2\ast & \ast & \ast & \ast & \ast \\
\lambda_{\alpha \beta \gamma \delta}\ast & c\ast & c\ast & c\ast & c\ast \\
\lambda_{\alpha \beta \gamma \delta}\ast & c\ast & c\ast & c\ast & c\ast \\
\end{bmatrix}.
\end{equation}
The determinant of \eqref{eq:my-det3} is $c(c^3 - \lambda_{\alpha \beta \gamma \delta})$ multiplied by a polynomial in $A^{(1)}, A^{(\alpha)}, A^{(\beta)}, A^{(\gamma)}, A^{(\delta)}$.
The most specialized case is $\alpha=\beta=\gamma$.  Computer verification with a random numerical instance proves the polynomial is not identically zero.
We deduce that $c^3 = \lambda_{\alpha \beta \gamma \delta}$.  
By symmetry, $c^3 = \lambda_{\pi(\alpha \beta \gamma \delta)}$ for all $\alpha, \beta, \gamma, \delta \in \{2, \dots, n\}$ with $\alpha, \delta$  distinct and all permutations $\pi$.
In other words, $\lambda$-entries with no $1$-indices and non-identical indices equal $c^3$.  

\medskip

Steps 1, 2 and 3 show that  $\lambda$ takes the announced form.  So, $\lambda$ is rank-$1$ off the diagonal.  
This finishes the ``only if" direction.
Overall, we have proven that the $5 \times 5$ minors of the $3n \times 27n^3$ flattenings of $\mathbf{Q}$ give algebraic relations on $\{Q^{(\alpha \beta \gamma \delta)} : \alpha, \beta, \gamma, \delta \in [n] \}$ with the desired \nolinebreak properties.

\end{document}
